\begin{lemma}[Impossibility of the negated candidate bound]\label{lem:step-015-negated-target}
Under Assumptions~\ref{assump:candidate-regime},
\ref{assump:central-dp}, and
\ref{assump:distribution-free-realizable-pac},
Proposition~\ref{prop:step-005-certificate}, and
Proposition~\ref{prop:step-014-fixed-contradiction}, fix any candidate in
the order \eqref{eq:candidate-order}--\eqref{eq:candidate-map}, with the
privacy pair chosen between the integer tuple and the learner. If

\[
n<ak\log_2^*N,
\tag{1}
\]

then the preceding propositions yield a contradiction. Hence no
such fixed candidate can satisfy (1).
\end{lemma}

\begin{proof}
The constants in \eqref{eq:calibration}, including \(a=b_*/16\), and
\(N_0\) were fixed by
Proposition~\ref{prop:step-005-certificate} before the present
candidate was chosen. Apply that proposition to the fixed candidate and the
strict hypothesis (1). It yields, at the exact integer budget

\[
M=\max\left\{8,\left\lceil\frac{4n}{k}\right\rceil\right\},
\]

the strict hard-regime and source-privacy conclusions

\[
N\ge N_*,\qquad 8\le M<b_*\log_2^*N,
\qquad 0<\varepsilon\le0.1,
\qquad 0<\delta<\frac{d_*}{M^2\log M}.
\tag{2}
\]

The same proposition retains, rather than replaces, the primitive
conjunction

\[
0<\delta\le\frac{1}{n\log(n+1)},
\qquad
0<\delta\le\frac{c_\delta}{M^2\log(M+1)}.
\tag{3}
\]

The fixed learner still satisfies its exact
Assumption~\ref{assump:central-dp} and
Assumption~\ref{assump:distribution-free-realizable-pac} premises; neither
is altered by the scalar specialization (2). Therefore the complete
premise package for
Proposition~\ref{prop:step-014-fixed-contradiction} holds at this identical
candidate. Applying that proposition, without reopening any of its internal
dependencies, gives one fixed product instance
\((P_{\boldsymbol Q^*},c_{\boldsymbol t^*})\) for which

\[
\mathbb E_{\substack{
 S\sim(P_{\boldsymbol Q^*}^{c_{\boldsymbol t^*}})^n\\
 \rho_A}}
R_{P_{\boldsymbol Q^*}}
\bigl(A_{\rho_A}(S),c_{\boldsymbol t^*}\bigr)>2^{-9}
\tag{4}
\]

and, for literally the same fixed distribution, target, size-\(n\) iid
sample law, randomized learner, internal-coin law, and population risk,

\[
\mathbb E_{\substack{
 S\sim(P_{\boldsymbol Q^*}^{c_{\boldsymbol t^*}})^n\\
 \rho_A}}
R_{P_{\boldsymbol Q^*}}
\bigl(A_{\rho_A}(S),c_{\boldsymbol t^*}\bigr)\le2^{-12}.
\tag{5}
\]

But

\[
2^{-9}=8\cdot2^{-12}>2^{-12},
\tag{6}
\]

so no real number can satisfy both (4) and (5). This contradicts (1).

There is no exceptional boundary branch. Proposition~\ref{prop:step-005-certificate} derives (2)--(3) with the same
strict inequalities when \(N=N_0\), \(M=8\), \(n=1\), \(n<k\), or
\(k=2\), and Proposition~\ref{prop:step-014-fixed-contradiction} applies to that same
candidate and exact fixed-instance experiment. No floor, ceiling, endpoint,
privacy parameter, or sample size is changed in this invocation. \end{proof}

\begin{proposition}[Exact nonasymptotic candidate closure]\label{prop:step-015-exact-closure}
Fix the universal constants in \eqref{eq:calibration} and the fixed
\(N_0\) before
candidate quantification. Under Assumptions~\ref{assump:candidate-regime},
\ref{assump:central-dp}, and
\ref{assump:distribution-free-realizable-pac}, every candidate quantified
in that same candidate order satisfies

\[
n\ge a k\log_2^*N.
\tag{7}
\]

This is a deterministic pointwise implication for the exact candidate. It
makes no assertion outside the approved candidate regime.
\end{proposition}

\begin{proof}
Fix an arbitrary candidate covered by the proposition. If (7) failed, then
the total order on the real numbers would give its exact strict negation

\[
n<ak\log_2^*N.
\tag{8}
\]

No integer rounding is inserted into this logical negation. But
Lemma~\ref{lem:step-015-negated-target} proves that (8) is impossible for a
candidate satisfying the three primitive assumptions. Therefore (7) holds.
Because the constants were fixed before the arbitrary candidate and the
candidate was otherwise arbitrary within that order, the conclusion has the
quantifier order stated in the formalized goal. \end{proof}
